% sec11_further.tex
% Section 11: further remarks and open problems.
% Labels defined here: sec:further, sec:onen-open, conj:onen, sec:tilt, conj:tilt, sec:nodes.

\section{Further remarks and open problems}\label{sec:further}

Theorem~\ref{thm:analytic} gives the pointwise law with rate $O(1/n)$ for every $p>0$, and Theorem~\ref{thm:onen}
identifies the $1/n$ term for pure powers with $p>1$. This section records what the computations of
Section~\ref{sec:numerics} suggest beyond these results, states two conjectures, and lists the directions in which the
method stops. Throughout, $a_n/(2b_n)$ pairs the diagonal entry at $P_n$ with the off-diagonal entry $b_n$ between
$P_{n-1}$ and $P_n$.

\subsection{\texorpdfstring{The $1/n$ term for $p\le1$}{The 1/n term for p <= 1}}\label{sec:onen-open}
The fits of Table~\ref{tab:onen} for $p=\frac12$ and $p=1$ agree with the values of Theorem~\ref{thm:onen} as well as
the error estimates allow.

\begin{conjecture}\label{conj:onen}
Let $0<p\le1$ and $s>1$. Then the pure powers satisfy
\[
  n^pb_n(\sigma_{p,s})=c(p)\Bigl(1-\frac{s-1}{2n}+o\Bigl(\frac1n\Bigr)\Bigr),\qquad
  \frac{a_n(\sigma_{p,s})}{2b_n(\sigma_{p,s})}=1-\frac{p}{2n}+o\Bigl(\frac1n\Bigr).
\]
\end{conjecture}

The evidence consists of the rows $p\le1$ of Table~\ref{tab:onen} ($s\in\{\frac32,2,3,4\}$, $n\le1600$) and the
columns $\ell\equiv1$ of Table~\ref{tab:numerics} at $p=\frac12$ and $p=1$. At $p=\frac12$ the remainder does not
appear to be $O(n^{-2})$: the fits need the powers $n^{-1-jp}$, and in Table~\ref{tab:numerics} the deviation from
$-\frac12$ decays like a term of order $n^{-1-p}$ in $n^pb_n/c(p)$. This is why the conjecture is stated with $o(1/n)$
only. For $p=1$ the averaged form of the statement for $b_n$ is proved in \cite{Paper3}: for every $s>1$,
\[
  \sum_{j=1}^n\log\frac{j\,b_j(\sigma_{1,s})}{c(1)}+\frac{s-1}2\log n
\]
converges as $n\to\infty$. This is what the conjecture predicts after summation, but it does not imply the pointwise
statement.

The method of Sections~\ref{sec:classical} and~\ref{sec:second} does not reach $p\le1$. It rests on two facts that hold
for $p>1$: the classical hard-edge parametrix is available for every $\alpha>-1$, and the band factor of
Section~\ref{sec:second} is constant up to $O(n^{-p})=o(1/n)$. At $p=1$ the classical construction fails
(Remark~\ref{rem:quotient-p1}), and the comparison argument of Section~\ref{sec:hardedge} bounds the hard-edge
contribution by $O(1/n)$ without computing its $1/n$ term; the band factor is not flat either, since it equals
$1+z/(6n)+O(n^{-2})$, although this defect alone enters the coefficients only at order $n^{-2}$. For $p<1$ the outer
scalar of Appendix~\ref{app:outer} gives no flatness to $o(1/n)$ near $0$, and the numerics show remainders of order
$n^{-1-p}$ that a proof would have to produce. What is missing is a hard-edge parametrix with matching error
$O(n^{-2})$ at $p=1$, and one adapted to the seams of Appendix~\ref{app:outer} for $p<1$.

\subsection{Slowly varying factors at order \texorpdfstring{$1/n$}{1/n}}\label{sec:tilt}
For an analytic factor $\ell$ the columns $\ell$ of Table~\ref{tab:numerics} drift away from the pure-power values. The
drift is described by the following conjecture.

\begin{conjecture}\label{conj:tilt}
Let $p>0$, $s>1$, and let $\ell$ be analytic of angle $\vartheta$ for some $\vartheta>0$, with
$\ell(k)=\exp\Lambda(\log k)$. Then
\[
  \frac{n^pb_n(\sigma_{p,s,\ell})}{c(p)}=1-\frac{s-1-\Lambda'(\log n)}{2n}+o\Bigl(\frac{\Lambda'(\log n)}{n}\Bigr)
  \qquad\text{whenever}\quad n\Lambda'(\log n)\to\infty.
\]
\end{conjecture}

The conjecture says that, at order $1/n$, a slowly varying factor acts like a change of the exponent $s$ by the local
logarithmic derivative $\Lambda'(\log n)$, that is, like a hard-edge Bessel order that drifts with $n$. In
Table~\ref{tab:numerics} the difference between the first two columns is $0.91$ to $0.99$ times
$\frac12\Lambda'(\log n)$, and this ratio increases with $n$ in all three geometries. In the thin-strip example of
Section~\ref{sec:numerics}, where $\Lambda'$ oscillates without a sign, the $1/n$ term oscillates with it. Within our
method the conjecture amounts to computing the main term of the comparison estimate of Section~\ref{sec:hardedge}, as
Conway and Deift do for their weight \cite[(5.5)]{CD}; this would need a Bessel parametrix with a drifting order.

\subsection{Other Bessel orders, other nodes, rougher weights}\label{sec:nodes}
For pure powers with $p>1$ the classical parametrix of Section~\ref{sec:classical} works for every Bessel order
$\alpha>-1$. For analytic factors $\ell$, and for pure powers with $p\le1$, the comparison method is confined to
$|\alpha|<1$, and Bessel orders $\alpha\ge1$ are reached only through the Christoffel shift
(Theorem~\ref{thm:shift}). A Riemann--Hilbert analysis of the weight itself for $\alpha\ge1$ would require a different
comparison frame (Remark~\ref{rem:alpha-ge-1}); it would be needed to compute the $1/n$ term directly in that range.

At $p=1$ one can also work in the symmetric variable $y=\pm x^{1/2}$ of \cite{WLXZ}, where $0$ is an interior
Fisher--Hartwig point of the band. The comparison method then covers only $1<s<2$ directly, and $1<s<3$ after folding;
we have not pursued this.

Our nodes are exact powers $k^{-p}$. For nodes $x_k=f(k)$ with $f'(\xi)\sim-pC\xi^{-p-1}$ the averaged law is known
\cite[Corollary~1.2]{Paper1}, but the pointwise law $n^pb_n\to c(p)C$ is open. One expects it to hold when $f$ is smooth
in a sense analogous to Definition~\ref{def:analytic}, and to fail under rough displacements of the nodes, as it does
for rough weights: bounded rough weights $k^{-s}e^{\epsilon_k}$, $|\epsilon_k|\le\varepsilon$, can break the pointwise
law although the averaged law holds for them, as is proved in \cite{Paper3}.
